% GENERATED FILE. Edit canonical lessons, then run npm run generate:books.
% canonical-source-hash: fnv1a64:3f1dfa7c80220c52
% canonical-lessons: ML-C516-kuli, ML-C516-poti, ML-C516-atayuka, ML-C516-tatayuka, ML-C516-tatassappetuka, ML-R516-review-509-516, ML-R516-review-final-516

\chapter{Blocked Roads}
\label{ch:ml-blocked-roads}
\hlchaptermodality{516}

\begin{chapteropening}
\textbf{By the end of this chapter:} \emph{I can say a pothole, dust, to get blocked, to block and to be obstructed.}

Complete the last lesson of chapter 516: I can say a pothole, dust, to get blocked, to block and to be obstructed.
\end{chapteropening}

\section[\texorpdfstring{\ml{കുഴി} (kuḻi)}{കുഴി (kuḻi)}]{\ml{കുഴി} (kuḻi) — a pothole; a pit}

\label{lesson:ML-C516-kuli}



\begin{quote}
Before the new one: say the Malayalam for to get tangled, then the Malayalam for to grope for.
\end{quote}

\subsection*{You'll want to know: \ml{കുഴി}}
\textbf{\ml{കുഴി}} — \emph{kuḻi} — \textquotedblleft{}a pothole; a pit\textquotedblright{}.

\begin{quote}
\ml{വഴിയിലെ} \ml{വലിയ} \ml{കുഴിയിൽ} \ml{ബൈക്ക്} \ml{വീണു}.
\end{quote}

\begin{quote}
\emph{vaḻiyile valiya kuḻiyil baikkŭ vīṇu.}
\end{quote}

\begin{quote}
The motorbike fell into a big pothole in the road.
\end{quote}

Plural: \ml{കുഴികൾ}, \emph{kuḻikaḷ}.

Any hole in the ground: a pit, a grave, a pothole.

\begin{cousinweb}
\textbf{\ml{കുഴി}} (\emph{kuḻi}) is native Dravidian, \textquotedblleft{}a hollow, a hole, a pit\textquotedblright{} (DEDR 1818).
\end{cousinweb}

\begin{grammarlens}[title={What you've built}]
One more word for getting around.
\end{grammarlens}

\subsection*{Your turn}
\begin{itemize}
\raggedright
  \item \emph{Say it:} \emph{kuḻi}
  \item \emph{Say it:} \emph{kuḻi}, once more
  \item \emph{Say it:} \emph{vaḻiyile valiya kuḻiyil baikkŭ vīṇu.} — The motorbike fell into a big pothole in the road.
  \item \emph{Recall:} say the Malayalam for to get tangled, then the Malayalam for to grope for, then say \emph{kuḻi} again
\end{itemize}

\begin{tcolorbox}[breakable,colback=teal!4,colframe=teal!35!black,title={Before you move on}]
What does \ml{കുഴി} mean? (\textquotedblleft{}A pothole; a pit\textquotedblright{}.) Say it once more.
\end{tcolorbox}
\section[\texorpdfstring{\ml{പൊടി} (poṭi)}{പൊടി (poṭi)}]{\ml{പൊടി} (poṭi) — dust; powder}

\label{lesson:ML-C516-poti}



\begin{quote}
Before the new one: say the Malayalam for to grope for, then the Malayalam for a pothole.
\end{quote}

\subsection*{You'll want to know: \ml{പൊടി}}
\textbf{\ml{പൊടി}} — \emph{poṭi} — \textquotedblleft{}dust; powder\textquotedblright{}.

\begin{quote}
\ml{വേനലിൽ} \ml{വഴിയിൽ} \ml{നിറയെ} \ml{പൊടിയാണ്}.
\end{quote}

\begin{quote}
\emph{vēnalil vaḻiyil niṟaye poṭiyāṇŭ.}
\end{quote}

\begin{quote}
In the dry season the road is full of dust.
\end{quote}

Also any powder: \ml{മുളകുപൊടി}, chilli powder.

\begin{cousinweb}
\textbf{\ml{പൊടി}} (\emph{poṭi}) is native Dravidian, \textquotedblleft{}dust, powder\textquotedblright{} (DEDR 4481).
\end{cousinweb}

\begin{grammarlens}[title={What you've built}]
One more word for getting around.
\end{grammarlens}

\subsection*{Your turn}
\begin{itemize}
\raggedright
  \item \emph{Say it:} \emph{poṭi}
  \item \emph{Say it:} \emph{poṭi}, once more
  \item \emph{Say it:} \emph{vēnalil vaḻiyil niṟaye poṭiyāṇŭ.} — In the dry season the road is full of dust.
  \item \emph{Recall:} say the Malayalam for to grope for, then the Malayalam for a pothole, then say \emph{poṭi} again
\end{itemize}

\begin{tcolorbox}[breakable,colback=teal!4,colframe=teal!35!black,title={Before you move on}]
What does \ml{പൊടി} mean? (\textquotedblleft{}Dust; powder\textquotedblright{}.) Say it once more.
\end{tcolorbox}
\section[\texorpdfstring{\ml{അടയുക} (aṭayuka)}{അടയുക (aṭayuka)}]{\ml{അടയുക} (aṭayuka) — to get blocked; to be shut}

\label{lesson:ML-C516-atayuka}



\begin{quote}
Before the new one: say the Malayalam for a pothole, then the Malayalam for dust.
\end{quote}

\subsection*{You'll want to know: \ml{അടയുക}}
\textbf{\ml{അടയുക}} — \emph{aṭayuka} — \textquotedblleft{}to get blocked; to be shut\textquotedblright{}.

\begin{quote}
\ml{ചവറ്} \ml{കാരണം} \ml{തോടുകൾ} \ml{അടഞ്ഞു}.
\end{quote}

\begin{quote}
\emph{cavaṟŭ kāraṇaṁ tōṭukaḷ aṭaññu.}
\end{quote}

\begin{quote}
Because of the litter, the streams got blocked.
\end{quote}

Past: \ml{അടഞ്ഞു}, \emph{aṭaññu}.

The intransitive partner of \ml{അടയ്ക്കുക}, to close: a road, a drain or a nose gets blocked.

\begin{cousinweb}
\textbf{\ml{അടയുക}} (\emph{aṭayuka}) is native Dravidian, \textquotedblleft{}to be shut, to be blocked up\textquotedblright{} (DEDR 83).
\end{cousinweb}

\begin{grammarlens}[title={What you've built}]
One more word for getting around.
\end{grammarlens}

\subsection*{Your turn}
\begin{itemize}
\raggedright
  \item \emph{Say it:} \emph{aṭayuka}
  \item \emph{Say it:} \emph{aṭayuka}, once more
  \item \emph{Say it:} \emph{cavaṟŭ kāraṇaṁ tōṭukaḷ aṭaññu.} — Because of the litter, the streams got blocked.
  \item \emph{Recall:} say the Malayalam for a pothole, then the Malayalam for dust, then say \emph{aṭayuka} again
\end{itemize}

\begin{tcolorbox}[breakable,colback=teal!4,colframe=teal!35!black,title={Before you move on}]
What does \ml{അടയുക} mean? (\textquotedblleft{}To get blocked; to be shut\textquotedblright{}.) Say it once more.
\end{tcolorbox}
\section[\texorpdfstring{\ml{തടയുക} (taṭayuka)}{തടയുക (taṭayuka)}]{\ml{തടയുക} (taṭayuka) — to block; to stop}

\label{lesson:ML-C516-tatayuka}



\begin{quote}
Before the new one: say the Malayalam for dust, then the Malayalam for to get blocked.
\end{quote}

\subsection*{You'll want to know: \ml{തടയുക}}
\textbf{\ml{തടയുക}} — \emph{taṭayuka} — \textquotedblleft{}to block; to stop\textquotedblright{}.

\begin{quote}
\ml{പോലീസ്} \ml{പാലത്തിൽ} \ml{വണ്ടികൾ} \ml{തടഞ്ഞു}.
\end{quote}

\begin{quote}
\emph{pōlīsŭ pālattil vaṇṭikaḷ taṭaññu.}
\end{quote}

\begin{quote}
The police stopped the vehicles at the bridge.
\end{quote}

Past: \ml{തടഞ്ഞു}, \emph{taṭaññu}.

What is stopped takes no ending or -\ml{എ}: \ml{വണ്ടികൾ} \ml{തടഞ്ഞു}.

\begin{cousinweb}
\textbf{\ml{തടയുക}} (\emph{taṭayuka}) is native Dravidian, \textquotedblleft{}to be obstructed; to stop\textquotedblright{} (DEDR 3031).
\end{cousinweb}

\begin{grammarlens}[title={What you've built}]
One more word for getting around.
\end{grammarlens}

\subsection*{Your turn}
\begin{itemize}
\raggedright
  \item \emph{Say it:} \emph{taṭayuka}
  \item \emph{Say it:} \emph{taṭayuka}, once more
  \item \emph{Say it:} \emph{pōlīsŭ pālattil vaṇṭikaḷ taṭaññu.} — The police stopped the vehicles at the bridge.
  \item \emph{Recall:} say the Malayalam for dust, then the Malayalam for to get blocked, then say \emph{taṭayuka} again
\end{itemize}

\begin{tcolorbox}[breakable,colback=teal!4,colframe=teal!35!black,title={Before you move on}]
What does \ml{തടയുക} mean? (\textquotedblleft{}To block; to stop\textquotedblright{}.) Say it once more.
\end{tcolorbox}
\section[\texorpdfstring{\ml{തടസ്സപ്പെടുക} (taṭassappeṭuka)}{തടസ്സപ്പെടുക (taṭassappeṭuka)}]{\ml{തടസ്സപ്പെടുക} (taṭassappeṭuka) — to be obstructed; to be held up}

\label{lesson:ML-C516-tatassappetuka}



\begin{quote}
Before the new one: say the Malayalam for to get blocked, then the Malayalam for to block.
\end{quote}

\subsection*{You'll want to know: \ml{തടസ്സപ്പെടുക}}
\textbf{\ml{തടസ്സപ്പെടുക}} — \emph{taṭassappeṭuka} — \textquotedblleft{}to be obstructed; to be held up\textquotedblright{}.

\begin{quote}
\ml{മരം} \ml{വീണ്} \ml{ഗതാഗതം} \ml{തടസ്സപ്പെട്ടു}.
\end{quote}

\begin{quote}
\emph{maraṁ vīṇŭ gatāgataṁ taṭassappeṭṭu.}
\end{quote}

\begin{quote}
A tree fell and the traffic was held up.
\end{quote}

Past: \ml{തടസ്സപ്പെട്ടു}, \emph{taṭassappeṭṭu}.

\ml{തടസ്സം}, an obstacle, plus -\ml{പ്പെടുക}, to undergo.

\begin{cousinweb}
\textbf{\ml{തടസ്സപ്പെടുക}} (\emph{taṭassappeṭuka}) is a native Dravidian verb built on \ml{തടസ്സം}, an obstruction (DEDR 3031), with -\ml{പ്പെടുക}.
\end{cousinweb}

\begin{grammarlens}[title={What you've built}]
One more word for getting around.
\end{grammarlens}

\subsection*{Your turn}
\begin{itemize}
\raggedright
  \item \emph{Say it:} \emph{taṭassappeṭuka}
  \item \emph{Say it:} \emph{taṭassappeṭuka}, once more
  \item \emph{Say it:} \emph{maraṁ vīṇŭ gatāgataṁ taṭassappeṭṭu.} — A tree fell and the traffic was held up.
  \item \emph{Recall:} say the Malayalam for to get blocked, then the Malayalam for to block, then say \emph{taṭassappeṭuka} again
\end{itemize}

\begin{tcolorbox}[breakable,colback=teal!4,colframe=teal!35!black,title={Before you move on}]
What does \ml{തടസ്സപ്പെടുക} mean? (\textquotedblleft{}To be obstructed; to be held up\textquotedblright{}.) Say it once more.
\end{tcolorbox}
\section[(dialogue)]{More words for getting around}

\label{lesson:ML-R516-review-509-516}



\begin{quote}
Say the Malayalam for temperature and humidity.
\end{quote}

\subsection*{Your turn}
Say these aloud:

\begin{itemize}
\raggedright
  \item temperature, humidity, sultry heat, a chill, to get hot
  \item fog, hail, to cloud over, to grow cold, to rumble
  \item a forecast, vigilance, a heatwave, a cyclone, to snap and fall
  \item to overflow its banks, waterlogging, to float up, to go under, to divert
  \item a relief camp, disaster relief, to cave in, to be uprooted, to be in distress
  \item a drought, to crack open, to get scorched, to be disconnected, to stop working
  \item breakage, compensation, a hardship, to get tangled, to grope for
  \item a pothole, dust, to get blocked, to block, to be obstructed
\end{itemize}

\begin{tcolorbox}[breakable,colback=teal!4,colframe=teal!35!black,title={Before you move on}]
Temperature? (\textbf{\ml{താപനില}}, \emph{tāpanila}.) Humidity? (\textbf{\ml{ആർദ്രത}}, \emph{ārdrata}.) Sultry heat? (\textbf{\ml{ഉഷ്ണം}}, \emph{uṣṇaṁ}.) A chill? (\textbf{\ml{കുളിർ}}, \emph{kuḷir}.) To get hot? (\textbf{\ml{ചൂടാകുക}}, \emph{cūṭākuka}.) Fog? (\textbf{\ml{മൂടൽമഞ്ഞ്}}, \emph{mūṭalmaññŭ}.) Hail? (\textbf{\ml{ആലിപ്പഴം}}, \emph{ālippaḻaṁ}.) To cloud over? (\textbf{\ml{മൂടിക്കെട്ടുക}}, \emph{mūṭikkeṭṭuka}.) To grow cold? (\textbf{\ml{തണുക്കുക}}, \emph{taṇukkuka}.) To rumble? (\textbf{\ml{ഇരമ്പുക}}, \emph{irampuka}.) A forecast? (\textbf{\ml{പ്രവചനം}}, \emph{pravacanaṁ}.) Vigilance? (\textbf{\ml{ജാഗ്രത}}, \emph{jāgrata}.) A heatwave? (\textbf{\ml{ഉഷ്ണതരംഗം}}, \emph{uṣṇataraṁgaṁ}.) A cyclone? (\textbf{\ml{ചുഴലിക്കാറ്റ്}}, \emph{cuḻalikkāṟṟŭ}.) To snap and fall? (\textbf{\ml{പൊട്ടിവീഴുക}}, \emph{poṭṭivīḻuka}.) To overflow its banks? (\textbf{\ml{കരകവിയുക}}, \emph{karakaviyuka}.) Waterlogging? (\textbf{\ml{വെള്ളക്കെട്ട്}}, \emph{veḷḷakkeṭṭŭ}.) To float up? (\textbf{\ml{പൊങ്ങുക}}, \emph{poṅṅuka}.) To go under? (\textbf{\ml{മുങ്ങിപ്പോകുക}}, \emph{muṅṅippōkuka}.) To divert? (\textbf{\ml{വഴിതിരിച്ചുവിടുക}}, \emph{vaḻitiriccuviṭuka}.) A relief camp? (\textbf{\ml{ക്യാമ്പ്}}, \emph{kyāmpŭ}.) Disaster relief? (\textbf{\ml{ദുരിതാശ്വാസം}}, \emph{duritāśvāsaṁ}.) To cave in? (\textbf{\ml{ഇടിയുക}}, \emph{iṭiyuka}.) To be uprooted? (\textbf{\ml{കടപുഴകുക}}, \emph{kaṭapuḻakuka}.) To be in distress? (\textbf{\ml{വലയുക}}, \emph{valayuka}.) A drought? (\textbf{\ml{വരൾച്ച}}, \emph{varaḷcca}.) To crack open? (\textbf{\ml{വിണ്ടുകീറുക}}, \emph{viṇṭukīṟuka}.) To get scorched? (\textbf{\ml{കരിയുക}}, \emph{kariyuka}.) To be disconnected? (\textbf{\ml{മുറിഞ്ഞുപോകുക}}, \emph{muṟiññupōkuka}.) To stop working? (\textbf{\ml{നിലയ്ക്കുക}}, \emph{nilaykkuka}.) Breakage? (\textbf{\ml{കേടുപാട്}}, \emph{kēṭupāṭŭ}.) Compensation? (\textbf{\ml{നഷ്ടപരിഹാരം}}, \emph{naṣṭaparihāraṁ}.) A hardship? (\textbf{\ml{ബുദ്ധിമുട്ട്}}, \emph{buddhimuṭṭŭ}.) To get tangled? (\textbf{\ml{കുരുങ്ങുക}}, \emph{kuruṅṅuka}.) To grope for? (\textbf{\ml{തപ്പുക}}, \emph{tappuka}.) A pothole? (\textbf{\ml{കുഴി}}, \emph{kuḻi}.) Dust? (\textbf{\ml{പൊടി}}, \emph{poṭi}.) To get blocked? (\textbf{\ml{അടയുക}}, \emph{aṭayuka}.) To block? (\textbf{\ml{തടയുക}}, \emph{taṭayuka}.) To be obstructed? (\textbf{\ml{തടസ്സപ്പെടുക}}, \emph{taṭassappeṭuka}.)
\end{tcolorbox}
\section[(dialogue)]{Blocked Roads — once more}

\label{lesson:ML-R516-review-final-516}



\begin{quote}
Say the Malayalam for a pothole and dust.
\end{quote}

\subsection*{Your turn}
\begin{itemize}
\raggedright
  \item \emph{Say it:} a pothole, dust, to get blocked, to block, to be obstructed
\end{itemize}

\begin{tcolorbox}[breakable,colback=teal!4,colframe=teal!35!black,title={Before you move on}]
A pothole? (\textbf{\ml{കുഴി}}, \emph{kuḻi}.) Dust? (\textbf{\ml{പൊടി}}, \emph{poṭi}.) To get blocked? (\textbf{\ml{അടയുക}}, \emph{aṭayuka}.) To block? (\textbf{\ml{തടയുക}}, \emph{taṭayuka}.) To be obstructed? (\textbf{\ml{തടസ്സപ്പെടുക}}, \emph{taṭassappeṭuka}.)
\end{tcolorbox}
